\documentclass[11pt]{article}

% ── arXiv standard packages ──────────────────────────────────────────────
\usepackage[utf8]{inputenc}
\usepackage[T1]{fontenc}
\usepackage{amsmath,amssymb,amsthm}
\usepackage{graphicx}
\usepackage{booktabs}
\usepackage{hyperref}
\usepackage{geometry}
\usepackage{authblk}
\usepackage{xcolor}
\usepackage{listings}

\geometry{margin=1in}
\hypersetup{
  colorlinks=true,
  linkcolor=blue!50!black,
  citecolor=blue!50!black,
  urlcolor=blue!50!black,
}

\lstset{
  basicstyle=\ttfamily\small,
  breaklines=true,
  frame=single,
  columns=fullflexible,
}

\title{\bfseries Staging Localization and Zero-Model Controls for\\
Connectome-Constrained Controllers:\\
Five Pre-Registered Negative Results}

\author[1]{Lee\thanks{Correspondence: \texttt{ljjohn@hotmail.com}}}
\affil[1]{Independent researcher}

\date{\today}

\begin{document}
\maketitle

\begin{abstract}
Connectome-constrained models of neural control are increasingly reported to
``produce behaviours''. We argue that two measurement practices are missing from
this literature, and we quantify what their absence costs. First, a widely used
separability criterion---the Fisher discriminant ratio $\mathrm{Sb}/\mathrm{Sw}$
thresholded at $1.0$---has a non-zero null expectation of $(K{-}1)/[K(N{-}1)]$;
we show the criterion rejects data that decodes at 92\% accuracy. Second, decoder
accuracies are routinely reported without a strong null control; we show that
random linear projections, sparse random projections, and random neuron subsets
all reach 100\% on the same data, so such numbers do not identify structure.
We then report five pre-registered negative results from a system that drives a
29-DoF humanoid with a frozen \emph{Drosophila} whole-brain connectome
(166{,}700 neurons, 25.58M synapses, LIF dynamics): descending-neuron and
motor-neuron populations carry $\approx 0$ bits of task-family information; a
nonlinear readout underperforms a linear one; and a frozen LIF network exhibits
no habituation. A staging measurement localizes the information loss to spike
transduction: information arrives at the descending neurons' synaptic current at
100\% and leaves their spike output at 82.7\%. We also report what survives: the
connectome's \emph{topology} carries causal timing information, with
degree-preserving rewiring reducing a measured 102.8\,N punch contact to 0\,N
while an Erd\H{o}s--R\'enyi graph of matched size fails to trigger at all.
We release the measurement toolkit and a 24-source survey annotated by interface
paradigm.
\end{abstract}

\section{Introduction}

The connectome-as-controller literature has an interpretation problem that its
practitioners have begun to name. A model is instantiated from a connectome,
given a body, and reported to produce locomotion, grooming, or
escape~\cite{flygm,neuromechfly}. Whether
such a report supports the claim that the connectome \emph{computes} the
behaviour depends entirely on how the interface between network and body was
constructed---and that is the variable least often reported.

This paper is a negative-result contribution with two methodological
prescriptions. We do not claim a better controller. We claim that two common
measurement practices yield non-identifiable conclusions, show what happens when
they are corrected, and report what remains standing.

\paragraph{Contributions.}
\begin{enumerate}
  \item A \emph{staging localization} procedure: measure separability at every
        stage of the pipeline with the same decoder, so that the drop-off
        identifies the information-loss layer rather than an overall verdict.
  \item A correction to a common separability criterion, with the null
        expectation derived and the failure demonstrated quantitatively.
  \item A demonstration that random linear projections match the reported
        accuracy of several structural representations, motivating a mandatory
        control.
  \item Five pre-registered negative results, each closed by a criterion
        registered before the measurement.
  \item One positive result with a three-model null comparison: connectome
        topology carries causal timing information.
\end{enumerate}

\section{System under test}

A whole-brain connectome of adult male \emph{Drosophila} (MaleCNS v1.0;
166{,}700 neurons, 25{,}582{,}938 directed synapses) is instantiated as a
leaky integrate-and-fire network in C++. It receives proprioceptive input
(29 joint angles) through a sensory encoder into the ventral nerve cord and its
descending-neuron (DN) output drives, through a hand-designed readout, a
Unitree G1 humanoid simulated in MuJoCo with 29 actuated degrees of freedom.
The connectome parameters are frozen bit-for-bit: resting potential, reset
potential, threshold, membrane and synaptic time constants, refractory period,
conduction delay, and the synaptic weight matrix.

All measurements below were produced by C++ programs with the numbers written to
versioned JSON artifacts; no result depends on a Python computation.

\section{Staging localization}
\label{sec:staging}

The question ``does the connectome work?'' is unactionable. We instead measure
separability at each stage using an identical decoder: leave-one-out
nearest-centroid over standardized features, with a label-permutation null.

Table~\ref{tab:staging} shows the result for a five-class posture-family task
(300 trials, 60 per family).

\begin{table}[h]
\centering
\begin{tabular}{@{}lrr@{}}
\toprule
Stage & Dimensionality & LOO accuracy \\
\midrule
S0 joint angles                & 29    & 100.0\% \\
S1 encoder drive               & 6370  & 100.0\% \\
S2 sensory neuron spikes       & 6370  & 100.0\% \\
S3a DN synaptic current        & 1314  & 100.0\% \quad $\leftarrow$ arrives intact \\
S3  DN spikes                  & 1314  & 82.7\% \quad $\leftarrow$ loss occurs here \\
\bottomrule
\end{tabular}
\caption{Per-stage separability. Information reaches the descending neurons'
synaptic current without loss and is destroyed by spike transduction.}
\label{tab:staging}
\end{table}

The actionable reading is not ``the connectome does not carry information'' but
``spike transduction is lossy''. These imply different literatures and different
remedies. Without staging, the project would have abandoned the wrong layer.

A ceiling scan over drive gain ($\times 1.6$ to $\times 12.8$), four time
windows within the 120-frame trial, and two front-end codes (rate vs.\ Gaussian
population) reaches 92.0\% and never attains 100\%. The residual loss is
attributable to the threshold-and-refractoriness nonlinearity rather than to
insufficient drive, consistent with analyses that treat spike generation as an
explicit information bottleneck~\cite{snrib}.

\section{Two measurement errors}
\label{sec:errors}

\subsection{The Fisher ratio threshold has a non-zero null}

We initially decided separability by requiring the Fisher discriminant ratio
$\mathrm{Sb}/\mathrm{Sw} > 1.0$, where $\mathrm{Sb}$ is the between-class
scatter and $\mathrm{Sw}$ the within-class scatter. The criterion reported
``inseparable'' and the project came close to abandoning its central line of
work.

Under the null hypothesis of no class structure, the expectation of
$\mathrm{Sb}/\mathrm{Sw}$ is not zero but
\begin{equation}
  \mathbb{E}\left[\frac{\mathrm{Sb}}{\mathrm{Sw}}\right]
  = \frac{K-1}{K(N-1)},
\end{equation}
which equals $0.068$ for $K=4$ classes and $N=12$ samples per class. A threshold
of $1.0$ therefore demands near-zero within-class variance rather than
merely exceeding the noise floor.

On our own data, $\mathrm{Sb}/\mathrm{Sw} = 0.409$, which is $6\times$ the null
expectation (one-way ANOVA: $F = 6.00$, $\mathrm{df} = (3,44)$, $p = 0.0016$).
The corresponding leave-one-out decoding accuracy is 92\%. The criterion labelled
data that decodes at 92\% as inseparable.

A second defect compounds the first: computing intra-class dispersion as the mean
Euclidean distance to the centroid rather than as a standard deviation makes the
denominator scale as $\sqrt{d}$. At $d = 26$ dimensions, white noise inflates the
ratio's denominator by roughly a factor of five.

The remedy is not a better threshold but a permutation null: shuffle labels,
recompute, and compare against the resulting 95th percentile.

\subsection{Accuracies without strong null controls are not identifiable}

At one point several representations decoded at 100\%. We then ran the control
that should have come first: replace the neural pipeline with a random linear map
of the same output width.

\begin{center}
\begin{tabular}{@{}lr@{}}
\toprule
Representation & LOO accuracy \\
\midrule
Random Gaussian projection, $6370 \to 1314$ & 100.0\% \\
Sparse random projection (2\% density)      & 100.0\% \\
Random neuron subset (1314 of 6370)         & 100.0\% \\
\bottomrule
\end{tabular}
\end{center}

If a structure-free linear map saturates the metric, the metric does not
identify structure. This retroactively invalidated every accuracy figure the
project had produced, and it is the reason the toolkit accompanying this paper
runs this control automatically.

\section{Five pre-registered negative results}
\label{sec:negative}

Each result below was closed by a criterion registered \emph{before} the
measurement. This ordering is the only reason the results are informative rather
than post-hoc.

\subsection{N-1: DN population carries $\approx 0$ bits of task information}

Three named DN groups were examined across five posture families. One group fires
2--3 spikes per 240\,ms window with 2.5\% inter-family variation; an all-zero-joint
condition and a 1.9\,rad pose differ by 0.02 spikes, so the activity is an
input-independent floor rather than a response. A second group is active
(13--14\,Hz) but constant across families. A third sits on the decision threshold.

A linear readout of the full 1314-dimensional DN vector does reach 82.7\%
leave-one-out accuracy, so the space is decodable in a narrow sense. The features
that action selection would use, however, are flat. Both statements hold; a report
of only the first turns a 0-bit result into an apparent 83\% success.
A plausible mechanism is that a small number of strong connections dominate
predictions while the remainder constrains
little~\cite{strongconn}, so group-averaged readouts dilute whatever the sparse
strong subset carries.

\subsection{N-2: a hand-designed 3-group binary code covers 2 classes, not 7}

With three binary group signatures, $\chi^2 = 17.35$ ($\mathrm{df} = 9$,
$p = 0.043$), so the code is above chance. Mutual information is
$0.313$ bits against a 2-bit ceiling, and per-trial codes overlap across all four
families (one code appears in every family: 3, 4, 8, and 6 times respectively).
Aggregated per family, exactly two classes separate.

\subsection{N-3: nonlinear readout underperforms linear}

A 256--128 ReLU MLP reaches $74.1\% \pm 1.9$; ridge regression reaches
$82.7\% \pm 0.0$. The difference is $-8.5\% \pm 1.9$ with a confidence interval
excluding zero. At 36{,}000 samples the nonlinear hypothesis is
over-parameterized. Controls behaved as required: random-weight MLP 16.7\%,
shuffled labels 18.9\%, majority class 20.0\%, against a chance level of 20\%.

\subsection{N-4: motor-neuron channel rates also carry $0$ bits}

The motor-neuron layer is not silent: approximately 1{,}500 spikes per trial with
about 300 of 708 neurons active, and 100\% of trials non-zero. Across families the
six-channel rates vary by 4\% and 3\% respectively. The cause is structural rather
than a channel-selection error: the drive is broad-spectrum, so flexor and
extensor pools co-fire and cancel, with $\sum \mathrm{dir}\cdot\mathrm{rate}
\approx 0.000$. Changing the channel table does not alter the information content
of a broad-spectrum upstream drive.

\subsection{N-5: a frozen LIF network has no habituation dynamics}

With identical input repeated ten times, the early-window response is 261 spikes
on every repetition: ratio $1.000$ against a pre-registered criterion of
$\leq 0.70$. Across 44 trials with zero jitter, responses are bit-identical.
A frozen LIF network with no plasticity, no noise, and no adaptation current is a
deterministic map, so no habituation mechanism exists. The direction was closed
under its own pre-registered stop condition rather than by tuning.

\section{What survives: topology carries timing information}
\label{sec:positive}

Against the negative results above, one claim is supported by a three-model null
comparison. In a closed loop where the readout triggers a punch, we substituted
the connectome with two null graphs of identical size and edge count.

\begin{table}[h]
\centering
\begin{tabular}{@{}llrr@{}}
\toprule
Brain & Trigger & Contact frames & Peak force \\
\midrule
Real connectome                    & @36\,ms & 9 & \textbf{102.8\,N} \\
Degree-preserving rewiring         & @88\,ms & 0 & 0.0\,N \\
Erd\H{o}s--R\'enyi (matched size)  & never   & 0 & 0.0\,N \\
\bottomrule
\end{tabular}
\caption{Three-model null comparison. Degree-preserving rewiring retains each
neuron's in- and out-degree exactly and changes only edge targets.}
\label{tab:null}
\end{table}

Degree-preserving rewiring alone destroys the behavior, reducing 102.8\,N to
0\,N. The contribution is therefore topological rather than an artifact of the
degree distribution. The Erd\H{o}s--R\'enyi graph fails to trigger at all,
indicating that both topology and degree structure are required.

This result is causal and falsifiable, and it is the claim we are prepared to
defend: a frozen connectome's specific wiring is sufficient to drive a physical
humanoid into a measurable punch, for a period determined by its internal
dynamics.

\section{Literature context}

A 24-source survey accompanying this paper annotates each reference by
\emph{interface paradigm}, because that variable determines what a result
supports:

\begin{center}
\begin{tabular}{@{}llp{6.3cm}@{}}
\toprule
& Paradigm & Who does the work \\
\midrule
P1 & connectome as information processor, no learning & the wiring alone \\
P2 & frozen connectome + learned interfaces           & the adapters \\
P3 & connectome as structural prior                   & a trained policy \\
P4 & learned interface absorbs everything             & the adapter; connectome decorative \\
\bottomrule
\end{tabular}
\end{center}

Reports of the form ``the connectome produces behaviours'' are almost always P1
with reflex-class behaviours, or P4 where the interface performed the work.
Neither supports a claim that the connectome computes.
The digital-sphinx result~\cite{digitalsphinx} is the clearest case: a
\emph{C.~elegans} connectome driving a \emph{Drosophila} body produced realistic
walking, with the interface supplying the behavior.
Whole-CNS emulation~\cite{wholecns} is titled, in part, ``behaviours that do not
emerge''. P2 is the only paradigm with a genuine positive
precedent~\cite{flycontrol}. Connectome constraints also predict measured neural
responses well when the task matches the circuit's evolved
function~\cite{flyvis}, which is precisely the condition our task violates: a fly
connectome asked to classify humanoid punch families., and even there the authors report a 38.6\% episode
failure rate against 7.0\% for an MLP baseline, so the prior trades robustness
for speed.

The system in this paper occupies a configuration absent from all 24 sources: a
contact-dense task, a frozen connectome, and \emph{hand-designed} interfaces. Our
finding is that such interfaces carry 0 bits of intent information while the
topology carries real timing information---a combination no prior work reports.

\section{Recommendations}

\begin{enumerate}
  \item \textbf{Localize before concluding.} Report separability per stage with a
        fixed decoder. A single end-to-end number cannot distinguish a lossy
        transduction layer from an uninformative representation.
  \item \textbf{Use permutation nulls, not thresholds.} Derive the null
        expectation for any statistic you threshold. For
        $\mathrm{Sb}/\mathrm{Sw}$ it is $(K{-}1)/[K(N{-}1)]$, not zero.
  \item \textbf{Always run a random-projection control at matched width.} If it
        matches, the metric does not identify structure.
  \item \textbf{Name the interface paradigm.} P4 results always produce behavior,
        which is precisely why they cannot support connectome claims.
  \item \textbf{Include a case that must fail.} A criterion that cannot reject is
        not a criterion. The released toolkit contains a deliberately
        non-separable input whose rejection verifies the test.
  \item \textbf{Register stop conditions.} All five negative results here were
        closed by conditions written before the measurement, which is what
        prevents the result from being tuned into a positive one.
\end{enumerate}

\section{Availability}

Code (MIT) and documentation (CC~BY~4.0):
\begin{center}
\url{https://github.com/ljjohn77/connectome-null-controls}
\end{center}
The package contains the staging probe with permutation null and
random-projection control, the five negative results with the number that closed
each, the three null models, and the annotated 24-source survey. It contains no
third-party data; reproduction requires the publicly available MaleCNS
connectome and MuJoCo models, whose licenses are listed in the repository.

\section*{Acknowledgements}

The author thanks the maintainers of the publicly released connectome datasets
and simulation models on which this work depends, and the authors of the
surveyed negative and critical results, whose willingness to publish
non-confirmatory findings made this analysis possible.

\begin{thebibliography}{99}

\bibitem{wholecns}
Embodied whole-CNS emulation of \emph{Drosophila melanogaster}: behaviours that
emerge from the connectome, and behaviours that do not.
\emph{arXiv:2610.33881}; dataset \emph{Zenodo} 23004436.

\bibitem{digitalsphinx}
The digital sphinx: can a worm brain control a fly body?
\emph{eLife Reviewed Preprint} 111516.

\bibitem{flycontrol}
FlyControl: a \emph{Drosophila} connectome prior enables adaptable embodied
control. \url{https://thu-neuroml.github.io/FlyControl/}

\bibitem{flygm}
Whole-brain connectomic graph model enables whole-body locomotion control in
fruit fly. \emph{arXiv:2602.17997}.

\bibitem{neuromechfly}
NeuroMechFly v2: simulating embodied sensorimotor control in adult
\emph{Drosophila}. \emph{Nature Methods} (2024).
doi:10.1038/s41592-024-02497-y.

\bibitem{flyvis}
Connectome-constrained networks predict neural activity across the fly visual
system. \emph{Nature} (2024).

\bibitem{strongconn}
Infrequent strong connections constrain connectomic predictions of neuronal
activity. \emph{Cell} (2025).

\bibitem{snrib}
Learning to time-decode in spiking neural networks through the information
bottleneck. \emph{NeurIPS} (2021).

\end{thebibliography}

\end{document}
